% GPT-6 Astra Pro; /root/fixed_surface; 22 November 2026 UTC.
% Effort not exposed. New proof after S22, initially for independent review.
\section[HD-F1-S23: the actual four-dimensional residual obstruction]{HD-F1-S23: an exact genus-five period certificate and the residual Atiyah map of the norm fourfold}
\label{sec:f1-norm-fourfold-period-obstruction}

There is a nonempty open family of the actual connected etale cyclic
triple covers of genus-five curves used in S13--S22 on which the Prym
period derivative has rank twelve. This is proved here by an explicit
nodal multiplication calculation and an actual smoothing of its finite
etale cover. For each cover in the resulting open family, the
characteristic action on the sixteen-dimensional polarized Weil
deformation space factors through a specified four-dimensional quotient:
\[
 \mathfrak D_P\longrightarrow(\ker\mu_C)^*
       \xrightarrow{\,\overline{\operatorname{ob}}_E\,}
       \ker\bigl(\sigma_E^2:\operatorname{Ext}^2_P(E,E)
                                      \longrightarrow HH_{-2}(P)\bigr),
 \qquad E=\mathcal O_Y.
 \tag{F1S23.1}
\]
All arrows are constructed below. In particular the characteristic
action on this sixteen-dimensional space has rank at most four.
The value of the residual map in (F1S23.1) has not been computed;
the statement supplies its exact source and factorization, and gives
no further assertion about the full degree-two trace.

\paragraph{1. A specified nodal curve and its actual triple cover.}
Let $C_0$ be the integral nodal curve obtained from $\mathbb P^1_t$
by identifying the following five disjoint pairs:
\[
 (a_i,b_i)_{i=1}^5=((0,1),(2,4),(3,8),(5,11),(7,13)).
 \tag{F1S23.2}
\]
It is stable and has arithmetic genus five. Fix
$\omega^2+\omega+1=0$, $\omega\ne1$. The line bundle
$\epsilon_0$ has trivial pullback to the normalization and gluing
multiplier $\omega$ from the fibre at $a_i$ to that at $b_i$,
at every node. Its third tensor power has the specified trivialization
induced by $\omega^3=1$. It is nontrivial: a global section of
its normalization pullback is constant, and the gluing equation
$c=\omega c$ forces that constant to vanish. Thus
\[
 \pi_0:Z_0=\operatorname{Spec}_{C_0}
   (\mathcal O_{C_0}\oplus\epsilon_0^{-1}
                         \oplus\epsilon_0^{-2})\longrightarrow C_0
 \tag{F1S23.3}
\]
is a finite etale degree-three cover, with multiplication given by
that trivialization. Its normalization is three copies of
$\mathbb P^1$, and the node gluing cyclically joins the three
sheets. Its dual graph is connected, so $Z_0$ is connected.
There are fifteen edges and three vertices; hence its arithmetic
genus is thirteen. Locally at each node the three branches of the
cover are three copies of the same nodal local ring, so the asserted
etaleness also follows directly at the nodes.

\paragraph{2. The exact canonical multiplication matrix.}
Put
\[
 Q(t)=\prod_{i=1}^5(t-a_i)(t-b_i),\qquad
 f_i(t)=\frac1{t-a_i}-\frac{\omega}{t-b_i},\quad
 g_i(t)=\frac1{t-a_i}-\frac{\omega^2}{t-b_i}.
 \tag{F1S23.4}
\]
Sections of $\omega_{C_0}\otimes\epsilon_0$ are the rational
differentials with simple poles at the ten branches, whose residues
at a pair are $(r_i,-\omega r_i)$. The residue theorem imposes
$(1-\omega)\sum_i r_i=0$. Thus the four differentials
$(f_i-f_5)dt$, $1\le i\le4$, form a basis. The conjugate four
$(g_j-g_5)dt$ give a basis for
$H^0(\omega_{C_0}\otimes\epsilon_0^{-1})$.
These are full section spaces: away from the prescribed poles
there are no additional differentials on $\mathbb P^1$.

A section of $\omega_{C_0}^{\otimes2}$ has the unique expression
$P(t)(dt)^2/Q(t)^2$ with $\deg P\le16$, subject to the five
node conditions
\[
       P(a_i)Q'(b_i)^2=P(b_i)Q'(a_i)^2\quad(1\le i\le5).
 \tag{F1S23.5}
\]
Indeed the generator of the dualizing line at a node is
$du/u=-dv/v$, whose square has equal leading coefficients on
the two branches. These are precisely (F1S23.5).
The five conditions are independent: evaluation at the ten
distinct points is surjective already on polynomials of degree at
most nine, and the five pairs of evaluation coordinates are
disjoint. Consequently this target has dimension $17-5=12$.

The actual multiplication map has polynomial columns
\[
 P_{ij}(t)=Q(t)^2(f_i-f_5)(g_j-g_5),\qquad1\le i,j\le4.
 \tag{F1S23.6}
\]
Their degrees are at most sixteen and they satisfy (F1S23.5),
as is also immediate from multiplying the two residue rules.
Take the $17\times16$ matrix $M$ whose row $r$ is the coefficient
of $t^r$, $0\le r\le16$, and whose columns are lexicographically
ordered by $(i-1,j-1)$. All entries lie in $\mathbb Z[\omega]$.
Reduce by the ring map
\[
       \mathbb Z[\omega]\longrightarrow\mathbb F_{19},
                           \qquad\omega\longmapsto7.
 \tag{F1S23.7}
\]
Here $7^2+7+1=0$ modulo nineteen, and all ten branch points
remain distinct. The minor on rows $0,1,\ldots,11$ and columns
$0,1,\ldots,10,15$ has determinant six in $\mathbb F_{19}$.
For a short arithmetic verification, ordinary elimination without
row exchanges gives the successive nonzero pivots
\[
             9,17,10,17,8,16,8,15,5,14,16,15,
 \qquad\text{whose product is }6\pmod {19}.
 \tag{F1S23.8}
\]
The defining polynomial formulas, the full reduced minor, all
eighty node-matching identities and this elimination are recorded
in \texttt{check\_norm\_period\_rank.py} and its exact certificate.
In particular the minor is nonzero in $\mathbb Z[\omega]$.
The target bound just proved now gives
\[
 \operatorname{rank}\bigl(
 H^0(\omega_{C_0}\epsilon_0)\otimes
 H^0(\omega_{C_0}\epsilon_0^{-1})
                 \longrightarrow H^0(\omega_{C_0}^{\otimes2})\bigr)=12.
 \tag{F1S23.9}
\]

\paragraph{3. Smoothing the specified cover and preserving the rank.}
Take a proper projective smoothing of the stable curve $C_0$ over
a complete local characteristic-zero base, with all five node
parameters nonzero on its smooth locus. Such smoothings are obtained
by the local equations $uv=t_i$: the obstruction group for a
nodal curve is $\operatorname{Ext}^2(\mathbb L_{C_0},\mathcal O_{C_0})=0$.
To see the vanishing, the local node has a length-one cotangent
resolution, its sheaf $\mathcal Ext^1$ has finite support, and
the remaining global degree-two cohomology of a coherent sheaf
on the curve vanishes. The resulting compatible proper formal
deformation is effective by the projective formal-existence theorem.

The actual finite etale cover (F1S23.3) extends uniquely to this
proper henselian family. This is the equivalence of finite etale
covers in Stacks, Lemma~58.9.1, Tag~0A48; its proof first lifts
over each infinitesimal thickening and then algebraizes the finite
algebra. The cyclic automorphisms and their group law extend by
full faithfulness. Equivalently, each lifted node is plumbed using
the same parameter as its image, with the deck transformation
permuting its three copies. Connectedness persists: upper
semicontinuity of $h^0(\mathcal O_{Z_s})$ and
$h^0(\mathcal O_{Z_0})=1$ force it to be one near the central
fibre. The resulting smooth fibres are the required connected
etale triples, with their actual eigensheaf $\epsilon_s$.

Relative duality and $h^0(\epsilon_0^{\pm1})=0$ give
$h^1(\omega_{C_0}\epsilon_0^{\pm1})=0$. Also
$h^1(\omega_{C_0}^{\otimes2})=0$, because its Serre dual is
$H^0(\omega_{C_0}^{-1})=0$. Hence the three section spaces in
(F1S23.9) are the fibres of vector bundles of ranks $4,4,12$,
with formation commuting with base change. Multiplication is
a morphism between those bundles. A nonzero maximal minor at
the central point remains a unit in the local base. It follows
that the same multiplication is onto on smooth fibres. The finite
cover and its multiplication morphism have finite presentation,
so they spread to an algebraic neighbourhood; smooth complex
fibres in this neighbourhood give a nonempty open locus.

Let $\mathcal R$ be a connected smooth Hurwitz component containing
one of these smooth covers, after choosing a fine level structure.
It is irreducible, and the surjectivity locus just obtained is
nonempty and open in $\mathcal R$. Intersect it with the dense
general-curve open of S13--S22. This gives the nonempty open family
used in the rest of this section. No assertion about a different
Hurwitz component is needed in this argument.

\paragraph{4. The actual period derivative and its four-dimensional quotient.}
For a cover $\pi:Z\to C$ in that open family, write
\[
 U=H^0(C,K_C\epsilon),\quad V=H^0(C,K_C\epsilon^{-1}),\qquad
 \mu_C:U\otimes V\longrightarrow H^0(C,K_C^2).
 \tag{F1S23.10}
\]
Both $U,V$ have dimension four, $\mu_C$ is onto, and
$\dim\ker\mu_C=4$. The actual cover determines its Prym
$P\subset J(Z)$ with restriction polarization $\beta$ and its
order-three action. Let $\mathfrak D_P\subset H^1(P,T_P)$ be
the tangent space preserving that polarization and action.

The nontrivial character part of $H^0(Z,K_Z)$ is $U\oplus V$,
via the eigensheaf description of $\pi_*K_Z$. The inclusion of
$P$ into $J(Z)$ identifies its Hodge structure with this nontrivial
part, using the actual projector
$1-(1+\sigma+\sigma^2)/3$. A polarization-preserving tangent
is symmetric for the polarization pairing. Commutation with
$\sigma$ leaves its two off-diagonal blocks, one the transpose
of the other. Thus there is an isomorphism
\[
 \jmath:\mathfrak D_P\xrightarrow{\sim}(U\otimes V)^*,
 \qquad
 \jmath(\kappa)(s\otimes t)
          =\langle\kappa(s),t\rangle_Z.
 \tag{F1S23.11}
\]
Here the bracket is Serre duality on $Z$, after applying the
specified nontrivial-character identification; this fixes the
scale of the displayed map.

The tangent space of the cover space is $H^1(C,T_C)$, since
finite etale covers lift uniquely across infinitesimal thickenings.
For $v$ in that space the variation of the upstairs holomorphic
forms is cup product with $\pi^*v$. Serre duality and the trace
of the degree-three map therefore give the exact identity
\[
 \jmath(d\operatorname{Pr}(v))(s\otimes t)
   =\langle\pi^*v,st\rangle_Z
   =\langle v,\operatorname{Tr}_{\pi}(st)\rangle_C
   =3\langle v,\mu_C(s\otimes t)\rangle_C.
 \tag{F1S23.12}
\]
The product of the opposite-character eigensections is the
pullback of $\mu_C(s\otimes t)$, which proves the last factor
three with the trace convention in (F1S23.11).
This is also the cotangent-period calculation of Lange--Ortega,
Proposition~4.1, with the trace scale made explicit here.
It proves injectivity of $d\operatorname{Pr}$, and the restriction
of $\jmath(\kappa)$ to the actual subspace $\ker\mu_C$ gives
the exact sequence
\[
 0\longrightarrow H^1(C,T_C)
  \xrightarrow{d\operatorname{Pr}}\mathfrak D_P
  \xrightarrow{\mathrm{res}\circ\jmath}(\ker\mu_C)^*
  \longrightarrow0.
 \tag{F1S23.13}
\]
Indeed the annihilator of $\ker\mu_C$ is
$\operatorname{im}\mu_C^*$, and (F1S23.12) identifies that
image with the actual period tangent. Thus the dimensions are
$12,16,4$ along exactly these maps.

\paragraph{5. The norm sheaf exists over those actual twelve directions.}
Over a sufficiently small smooth base in the open family, label a
component of the relative canonical norm fibre and choose a local
origin in its Prym torsor. S13 gives a proper smooth family
$\mathcal T\to B$ and an actual relative Abel morphism
$\mathcal a:\mathcal T\to\mathcal P$. Put
\[
                 \mathcal E=R\mathcal a_*\mathcal O_{\mathcal T}.
 \tag{F1S23.14}
\]
Both total spaces are smooth after taking this smooth base.
Properness gives a bounded coherent direct image, and regularity
of the total ambient space makes it perfect. Derived base change
identifies its derived fibre with
$Ra_{s*}\mathcal O_{T_s}$. The base-change square is
Tor-independent because both source and target are flat over
$B$; the same conclusion follows from a flat relative resolution.
S16--S17 identify every such fibre with the degree-zero sheaf
$\mathcal O_{Y_s}$. The relative Tor-amplitude criterion then
shows that $\mathcal E$ is a sheaf flat over $B$.
In particular (F1S23.14) gives an actual perfect extension of
$E=\mathcal O_Y$ over every first-order cover deformation.

For any $\kappa\in H^1(P,T_P)$ define
\[
 \operatorname{ob}_E(\kappa):
 E\xrightarrow{\operatorname{At}(E)}E\otimes\Omega_P^1[1]
  \xrightarrow{1_E\otimes\kappa}E[2].
 \tag{F1S23.15}
\]
This is the characteristic Hochschild action of the
$H^1(T_P)$ summand, with the diagonal Atiyah-class convention.
The actual extension (F1S23.14) makes (F1S23.15) vanish on
$\operatorname{im}d\operatorname{Pr}$. One can read this
directly from the diagonal adjunction extension: a lifting
perfect complex splits its obstruction two-extension.
The precise published comparison is Huybrechts--Thomas,
Theorem~2.10 and Lemma~3.2, or Corollary~3.4, in the corrected
2013 author version. The universal diagonal morphism there
is exactly the two-arrow composition in (F1S23.15).

By (F1S23.13), the restriction of this actual map to
$\mathfrak D_P$ therefore descends uniquely to a map
\[
          \overline{\operatorname{ob}}_E:
             (\ker\mu_C)^*\longrightarrow\operatorname{Ext}^2_P(E,E),
 \qquad
 \overline{\operatorname{ob}}_E(\jmath(\kappa)|_{\ker\mu_C})
                          =\operatorname{ob}_E(\kappa).
 \tag{F1S23.16}
\]
Independence of the representative follows from the proved
vanishing on the kernel of the quotient, not from a choice of
a complement to the period tangent.

\paragraph{6. The exact trace-kernel target.}
In the $\beta$-unitary character basis $a_i,b_j$ of S18,
the sixteen polarized Weil tangents are spanned by
\[
   \kappa_{ij}=\frac\partial{\partial a_i}\otimes\bar b_j
              +\frac\partial{\partial b_j}\otimes\bar a_i,
                         \qquad1\le i,j\le4.
 \tag{F1S23.17}
\]
Their action on $\beta=\sqrt{-1}\sum(a_i\bar a_i+b_j\bar b_j)$
is zero: the two resulting anti-holomorphic two-forms cancel.
They also annihilate
$w=\lambda A\bar B+\bar\lambda B\bar A$ from S18.
In the first term, replacing any $a_i$ by $\bar b_j$ repeats
a factor of $\bar B$; the second term is the same calculation
with the two character spaces exchanged. Thus each tangent
annihilates $w$ and every power of $\beta$.
The entire character in S14 is a linear combination of these
classes, so
\[
                \kappa\mathbin{\lrcorner}\operatorname{ch}(E)=0
                        \quad\text{for every }\kappa\in\mathfrak D_P.
 \tag{F1S23.18}
\]
The full Hochschild trace-characteristic identity, with the
same diagonal action used in (F1S23.15), is
\[
 \sigma_E^2\bigl(\operatorname{ob}_E(\kappa)\bigr)
                    =\kappa\mathbin{\lrcorner}\operatorname{ch}(E).
 \tag{F1S23.19}
\]
This is the Buchweitz--Flenner identity already checked in S18.
Equations (F1S23.18)--(F1S23.19) prove the target assertion in
(F1S23.1). They retain the original sheaf, polarization,
character action and full trace. The exact four-dimensional
map (F1S23.16) remains to be evaluated; the construction above
does not assign it rank zero or positive rank.

\paragraph{7. The entire remaining self-Ext summand and its exact quotient trace.}
Write $I=I_{Y/P}$ and let
\[
 I\longrightarrow\mathcal O_P\xrightarrow{u}E
                   \xrightarrow{\partial}I[1]
 \tag{F1S23.20}
\]
be the actual ideal triangle. For every $j$, pushforward of
multiplication in $\mathcal O_Y$ defines
$s_j:H^j(Y,\mathcal O_Y)\to\operatorname{Ext}^j_P(E,E)$.
Precomposition with $u$ gives
$r_j:\operatorname{Ext}^j_P(E,E)\to H^j(Y,\mathcal O_Y)$,
and $r_js_j=1$: after the adjunction
$Li^*\mathcal O_P=\mathcal O_Y$, it is just the unit identity
for multiplication on $\mathcal O_Y$. This is the same scalar
section as S18. All neighbouring connecting maps from
$H^j(Y,\mathcal O_Y)$ in the Hom sequence of (F1S23.20)
are therefore zero. In degrees one and two that sequence gives
the actual split isomorphisms
\[
 \begin{split}
 \operatorname{Ext}^1_P(E,E)&\simeq
       H^1(Y,\mathcal O_Y)\oplus\operatorname{Hom}_P(I,E),\\
 \operatorname{Ext}^2_P(E,E)&\simeq
       H^2(Y,\mathcal O_Y)\oplus\operatorname{Ext}^1_P(I,E).
 \end{split}
 \tag{F1S23.21}
\]
The second summand map in degree $j$ is explicitly
$d_j(z)=z[1]\partial$. No local complete-intersection
condition on $Y$ is required for this ideal triangle.
By S21--S22 the two scalar summands have dimensions eight
and twenty-eight respectively.

Naturality of the Atiyah class at $u$ and
$\operatorname{At}(\mathcal O_P)=0$ imply
$\operatorname{ob}_E(\kappa)u=0$. Thus (F1S23.16)
has a unique lift through the injection $d_2$ to
\[
 \overline h_E:(\ker\mu_C)^*\longrightarrow
                              \operatorname{Ext}^1_P(I,E),
 \qquad d_2\overline h_E=\overline{\operatorname{ob}}_E.
 \tag{F1S23.22}
\]
This describes the actual embedded-obstruction contribution.
It can also be seen directly on deformations. Any perfect
first-order extension of $E$ is an ordinary sheaf flat over
the dual-number base, by its relative Tor-amplitude zero.
The obstruction to lifting the section $1\in H^0(Y,\mathcal O_Y)$
is a class in $H^1(Y,\mathcal O_Y)$. Tensoring with an
infinitesimal ambient line bundle changes this class by its
restriction from $H^1(P,\mathcal O_P)$; the Cech formula is
addition of the line transition cocycle to the differences of
local lifts of $1$. The restriction is an isomorphism by S21,
so an actual such line twist makes $1$ lift. Nakayama then
makes the resulting map from the deformed ambient structure
sheaf surjective. Its kernel is an ideal, giving an actual
flat embedded deformation of $Y$. Conversely a flat embedded
deformation supplies the structure-sheaf extension; it is
perfect locally by lifting a finite free resolution on the
smooth central ambient fibre and applying the flatness
criterion to its last syzygy. Thus the first-order object and
embedded deformation problems are compared by these actual
line twists and quotient maps.

Finally define the explicitly known twenty-eight-dimensional
subspace
\[
 \Sigma=\{\alpha\operatorname{ch}(E):
                    \alpha\in H^2(P,\mathcal O_P)\}\subset HH_{-2}(P).
 \tag{F1S23.23}
\]
S22 and the scalar trace identity say that
$\sigma_E^2s_2$ is an isomorphism from
$H^2(Y,\mathcal O_Y)$ onto $\Sigma$. Consequently the entire
remaining trace calculation has the exact form
\[
 \tau_Y:\operatorname{Ext}^1_P(I,E)
       \xrightarrow{d_2}\operatorname{Ext}^2_P(E,E)
       \xrightarrow{\sigma_E^2}HH_{-2}(P)
       \longrightarrow HH_{-2}(P)/\Sigma,
 \qquad \ker\sigma_E^2\simeq\ker\tau_Y.
 \tag{F1S23.24}
\]
The isomorphism sends a trace-zero self-Ext class to its
second component in (F1S23.21). Its inverse, for
$z\in\ker\tau_Y$, is the actual expression
\[
       d_2(z)-s_2\bigl((\sigma_E^2s_2)^{-1}
                                      \sigma_E^2d_2(z)\bigr).
 \tag{F1S23.25}
\]
The inverse in this formula is taken only on $\Sigma$;
membership follows from $\tau_Y(z)=0$. Equations
(F1S23.22)--(F1S23.24) give
$\operatorname{im}\overline h_E\subset\ker\tau_Y$.
The values of $\tau_Y$ on its full actual Ext source and
of $\overline h_E$ on its four-dimensional source remain
uncomputed. These exact morphisms are the retained gate.

\paragraph{Primary inputs and source scope.}
The period derivative is the actual cup-product and Serre-trace
map of Lange--Ortega, \emph{Prym varieties of cyclic coverings},
Geom.\ Dedicata~150 (2011), 391--403,
DOI~10.1007/s10711-010-9512-9,
Section~4, Proposition~4.1 and its full supplied proof.
Its Section~5 criterion requiring Clifford index at least five
does not apply to genus five; that criterion is not used.
The rank computation here is instead the explicit actual curve
matrix (F1S23.6)--(F1S23.9).
Finite etale lifting uses the full proof of Stacks
Lemma~58.9.1, Tag~0A48, together with the infinitesimal
equivalence in Lemma~58.8.3, Tag~0BQB.
The obstruction comparison is Huybrechts--Thomas,
\emph{Deformation-obstruction theory for complexes via Atiyah
and Kodaira--Spencer classes}, Math. Ann.~346 (2010),
545--569, using corrected arXiv:0805.3527v2 (2013):
Theorem~2.10, Lemmas~3.1--3.2, and Corollary~3.4, with
the complete supplied proof of Theorem~3.3 also read.
The relative object in (F1S23.14) provides the actual extension
needed here. Neither local tangent-cone rigidity nor a
classification of all deformations of $Y$ is used.
